% Standalone Overleaf entrypoint. Compile with pdfLaTeX; no external figures/bib file required.
\documentclass[10pt,twocolumn,letterpaper]{article}
\usepackage[margin=0.68in]{geometry}
\usepackage[T1]{fontenc}
\usepackage{lmodern}
\usepackage{microtype}
\usepackage{booktabs,tabularx,enumitem,balance}
\usepackage[hidelinks]{hyperref}
\usepackage{titlesec}
\setlength{\columnsep}{0.25in}
\setlength{\parindent}{1em}
\setlength{\parskip}{2pt}
\setlist[itemize]{leftmargin=*,nosep,topsep=3pt}
\titleformat{\section}{\normalfont\normalsize\bfseries}{\Roman{section}.}{0.5em}{}
\titlespacing*{\section}{0pt}{7pt}{3pt}
\title{\vspace{-1.5em}\Large\bfseries MedTrace Guard: Pharmaceutical Provenance with\newline Accepted Custody and Safety Gates}
\author{\normalsize Ananya Mishra, Hemanth Reddy Narendrababu\\\normalsize Sai Harsha Yerapothu, Ridham Shah\\\small Arizona State University\quad |\quad CSE 540: Engineering Blockchain Applications}
\date{\vspace{-1em}Fall B 2026\vspace{-1em}}
\begin{document}
\maketitle
\begin{abstract}
MedTrace Guard will prototype pharmaceutical batch tracking using Solidity smart contracts and a React dashboard. Beyond recording a custody timeline, the system will require recipient acceptance of each handoff and enforce recall, expiry, quantity, and reported-condition policies before distribution or dispensing. An out-of-range simulated temperature observation will place a persistent hold on a batch until a regulator records a review. Certificates will remain off-chain, with explicit SHA-256 commitments for byte-level verification. We will compare the prototype with an independently implemented simplified tracking baseline, measuring gas costs, unsafe-action rejection, and local read performance. The contribution is a focused, testable integration of acknowledgment and safety controls, rather than a claim of new cryptography or guaranteed drug authenticity.
\end{abstract}

\section{Problem and Motivation}
Pharmaceutical distribution involves organizations with separate responsibilities and records. A useful tracking system must identify who submitted a record, preserve custody history, and prevent actions inconsistent with recorded safety conditions. WHO's 2017 analysis estimated that about one in ten medical products in low- and middle-income countries was substandard or falsified~\cite{who}; this motivates traceability but does not establish a current universal prevalence rate.

The FDA describes DSCSA as enabling interoperable electronic tracing of certain prescription drugs at the package level~\cite{fda}. Our classroom system tracks batches and is \emph{inspired by} traceability needs; it will not claim DSCSA compliance, clinical validation, or the ability to authenticate physical medicine from a QR label.

\section{Related Work and Project Contribution}
Sylim et al. proposed blockchain-based pharmaceutical distribution tracking~\cite{sylim}. Consequently, combining pharmaceuticals and blockchain is not itself novel. Standard smart contracts, role controls, and content-addressed storage are also established tools. Our course-level contribution will be a coherent enforcement workflow and reproducible comparison:
\begin{itemize}
\item \textbf{Accepted custody:} a proposed shipment changes custody only when the named recipient accepts a live, nonce-bound request.
\item \textbf{Persistent safety holds:} reported temperature excursions block actions; later normal readings cannot silently clear the hold.
\item \textbf{Explainable verification:} public lookup separates registered provenance, policy eligibility, document integrity, and unavailable evidence.
\end{itemize}
Partial dispensing will preserve quantity accounting. These features extend the original descriptive tracking concept; they are not asserted to be first in the research literature.

\section{Stakeholders and Architecture}
The manufacturer registers batches; a distributor accepts and forwards custody; a pharmacy accepts and dispenses quantities; a regulator records recalls and reviews holds. A separate administrator manages approved wallet roles. Public viewers can inspect nonsensitive records without a wallet.

\begin{center}
\fbox{\parbox{0.91\columnwidth}{\centering
React dashboard\\[-1pt]
$\downarrow$ read RPC \qquad $\downarrow$ wallet-signed writes\\[-1pt]
Solidity registry, policy checks, and event history\\[3pt]
Off-chain certificate files / IPFS $\longrightarrow$\\[-1pt]
browser SHA-256 comparison with on-chain commitment}}
\end{center}

The planned stack is Solidity 0.8.30, ethers.js, React, Vite, and a local Ethereum-compatible development chain. The contract will be demonstrated locally; public-testnet deployment will use Polygon Amoy or Sepolia as agreed with the instructor. A small custom role registry will provide explicit permissions. No cryptocurrency payments or tokens are required.

\section{Data and Smart Contract Workflow}
Each batch will contain a generated ID, manufacturer batch code, medicine name, manufacturer and current custodian, production/expiry timestamps, initial/remaining units, permitted temperature range, lifecycle stage, safety flags, and certificate locator/digest. Manufacturer-plus-code duplicates will be rejected. Original metadata and history entries will have no edit/delete functions.

The lifecycle is manufacturer registration, distributor handoff/acceptance, pharmacy handoff/acceptance, and partial or complete dispensing. Handoffs will have a named recipient, nonce, and deadline. Role and custody checks will reject unauthorized callers, invalid routes, replayed acceptances, and invalid quantities. Recalls will cancel pending shipments and permanently block distribution. Expiry will be evaluated against the chain timestamp. A condition hold will persist until an authorized, evidence-backed review, without overriding recall or expiry.

Certificate files will be stored off-chain; the demo will support local synthetic files and actual IPFS locators. SHA-256 will be calculated over exact bytes, separately from the CID. IPFS content requires pinning for availability~\cite{ipfs}. Simulated temperature records will retain their source payload and commit its digest on-chain; this does not establish sensor truth.

\balance
\section{Scope, Evaluation, and Feasibility}
The prototype will support one manufacturer--distributor--pharmacy route per batch, with no splitting, merging, payments, real patient data, or regulatory integrations. The UI will provide registration/actions, paged batch and history views, QR links, public policy lookup, and certificate-byte checks.

An independently written \emph{simplified baseline} will record registration, direct custody transfers, recall, and quantity-limited dispensing. It is not a reproduction of an unavailable original implementation. Evaluation will include:
\begin{itemize}
\item successful lifecycle and adversarial tests for roles, nonces, expiry, recall, holds, and quantities;
\item gas receipts for registration and complete handoffs, counting both Guard transactions;
\item 10/50/100-batch local growth checks and paged-read timings;
\item certificate tampering and unavailable-document cases.
\end{itemize}
The baseline has less metadata/history; gas differences will be reported as total design costs, not isolated per-feature effects. Local measurements will not be presented as public-chain throughput or proof of production scalability. No experimental results are claimed in this proposal.

\section{Milestones and Limitations}
\begin{center}\small
\begin{tabularx}{\columnwidth}{@{}lX@{}}
\toprule
\textbf{2026 date} & \textbf{Planned deliverable}\\\midrule
Oct. 25 & Proposal and completed team contract\\
Nov. 1 & Contract draft, tests, GitHub README\\
Nov. 15 & Integrated prototype and 5-minute demo\\
Dec. 4 & Code, evaluation, 7--10 minute demo, 4--8 page IEEE report\\\bottomrule
\end{tabularx}\end{center}
The team will divide contracts/testing, UI/integration, off-chain records/deployment, and evaluation/reporting. Team names, allocations, and meeting agreements will be finalized in the separate contract.

Blockchain preserves recorded assertions; it cannot prevent dishonest inputs, copied QR labels, omitted sensor reports, or compromised keys. Public metadata, trusted onboarding, growing history, transaction costs, and storage availability remain limitations. A centralized system with signed audit logs is a valid alternative; the final report will compare its simpler governance and cost with the shared ledger's independent auditability.

\begin{thebibliography}{4}\footnotesize
\bibitem{who} World Health Organization, ``1 in 10 medical products in developing countries is substandard or falsified,'' Nov. 28, 2017. \href{https://www.who.int/en/news/item/28-11-2017-1-in-10-medical-products-in-developing-countries-is-substandard-or-falsified}{WHO news release}.
\bibitem{fda} U.S. FDA, ``Drug Supply Chain Security Act (DSCSA).'' \href{https://www.fda.gov/drugs/drug-supply-chain-integrity/drug-supply-chain-security-act-dscsa}{FDA overview}, accessed Oct. 23, 2026.
\bibitem{sylim} P. Sylim, F. Liu, A. Marcelo, and P. Fontelo, ``Blockchain Technology for Detecting Falsified and Substandard Drugs in Distribution: Pharmaceutical Supply Chain Intervention,'' \emph{JMIR Res. Protoc.}, vol. 7, no. 9, e10163, 2018, doi: \href{https://doi.org/10.2196/10163}{10.2196/10163}.
\bibitem{ipfs} IPFS Docs, ``Persistence, permanence, and pinning.'' \url{https://docs.ipfs.tech/concepts/persistence/}, accessed Oct. 23, 2026.
\end{thebibliography}
\end{document}
